\itemhead{2.1 案例 · 案例1（教材p.33）}{XJ-B1-C02-CA-P033-Y0234}
\label{XJ-B1-C02-CA-P033-Y0234}
{\footnotesize 湘教版教材p.33 / PDF第40页 / 2.1 案例 / 案例1}
\field{原题}{2002 年在北京召开的第 24 届国际数学家大会的会标如图 2.1-2（1）所示，其中的几何图形取材于我国古代数学家对勾股定理的证明，如图 2.1-2（2）。图中的正方形 $ABCD$ 由 4 个全等的直角三角形和小正方形 $EFGH$ 组成。设直角三角形的两条直角边长分别为 $a,b\ (a\ne b)$，那么正方形 $ABCD$ 的边长为 $\sqrt{a^2+b^2}$。
\bankfig{C02-CA033}{0.62}
又 4 个直角三角形的面积之和为 $4\times\frac12ab=2ab$，正方形 $ABCD$ 的面积为 $a^2+b^2$，并且正方形 $ABCD$ 的面积大于 4 个直角三角形的面积之和，于是可得不等式 $a^2+b^2>2ab\ (a\ne b)$。}
\field{参考答案}{由面积分割，$a^2+b^2=2ab+(a-b)^2>2ab$，其中 $a\ne b$ 保证严格大于；若允许 $a=b$，应改为 $\ge$。}
{\small \field{关联知识}{I03 基本不等式、等号与最值建模（教材帮p.52、53、54 / PDF61、62、63）}
\field{方法与易错点}{列正数条件与固定量，构造基本不等式，解等号条件并代回实际限制。 易错：把一个下界误当最小值；多个等号条件不能同时实现；漏单位和整数限制。}
\field{答案依据}{答案册未提供本栏目的独立对应答案；依据原题、教材概念补写参考回应。开放性总结不设唯一数值答案。}}
